\documentclass[11pt,a4paper]{article}
\usepackage[T1]{fontenc}
\usepackage{lmodern,microtype}
\usepackage[margin=27mm,headheight=14pt]{geometry}
\usepackage{amsmath,amssymb,amsthm,mathtools}
\usepackage{booktabs,longtable,array}
\usepackage[dvipsnames]{xcolor}
\usepackage{fancyhdr}
\usepackage[colorlinks=true,linkcolor=MidnightBlue,citecolor=MidnightBlue,urlcolor=MidnightBlue,breaklinks=true]{hyperref}
\usepackage{bookmark}
\usepackage{enumitem}
\setlist{itemsep=3pt,topsep=5pt}
\pagestyle{fancy}\fancyhf{}
\lhead{\small Unequal-twist Stieltjes jets}\rhead{\small ProveIt continuation}
\cfoot{\thepage}
\renewcommand{\headrulewidth}{0.3pt}
\numberwithin{equation}{section}
\newtheorem{theorem}{Theorem}[section]
\newtheorem{proposition}[theorem]{Proposition}
\newtheorem{corollary}[theorem]{Corollary}
\newtheorem{lemma}[theorem]{Lemma}
\theoremstyle{definition}\newtheorem{definition}[theorem]{Definition}
\theoremstyle{remark}\newtheorem{remark}[theorem]{Remark}
\DeclareMathOperator{\FP}{FP}
\DeclareMathOperator{\Res}{Res}
\DeclareMathOperator{\Li}{Li}
\DeclareMathOperator{\sgn}{sgn}
\newcommand{\dd}{\,\mathrm d}
\newcommand{\Z}{\mathbb Z}\newcommand{\N}{\mathbb N}\newcommand{\C}{\mathbb C}
\newcommand{\T}{\mathbb R/\mathbb Z}
\newcommand{\eps}{\varepsilon}
\newcommand{\wt}[1]{\widehat{#1}}
\newcommand{\rv}{\boldsymbol r}\newcommand{\mv}{\boldsymbol m}
\newcommand{\av}{\boldsymbol\alpha}\newcommand{\dv}{\boldsymbol d}
\newcommand{\tv}{\boldsymbol\theta}\newcommand{\qv}{\boldsymbol q}
\newcommand{\wv}{\boldsymbol w}
\newcommand{\stir}[2]{\genfrac{[}{]}{0pt}{}{#1}{#2}}
\newcommand{\pin}{\texttt{b1df802799851c62a860570c3ac78ee31b9f5c04}}
\newcommand{\pathref}[1]{\nolinkurl{#1}}
\allowdisplaybreaks[2]
\hypersetup{pdftitle={Unequal-Twist Stieltjes Jets: Monodromy Obstructions, Convergent Harmonic Expansions, and Gamma Convolutions},pdfauthor={Prepared for Vladimir Reshetnikov with ChatGPT},pdfsubject={Exact polylogarithm, Hurwitz-zeta, Gamma and generalized Stieltjes identities}}
\title{\textbf{Unequal-Twist Stieltjes Jets}\\[5pt]
\large Monodromy Obstructions, Convergent Harmonic Expansions,\\
Gamma Convolutions, and Multishift Zeta Finite Parts}
\author{Research continuation prepared for Vladimir Reshetnikov\\
\small ProveIt programme\quad\textbullet\quad Prepared with ChatGPT}
\date{10 October 2026}
\begin{document}
\maketitle
\begin{abstract}
The current ProveIt manuscript proves finite partial-fraction formulas for
integer-order kernels with unequal frequency shifts, but leaves mixed
Stieltjes order jets in a prescribed finite Hurwitz basket as a separate
question. We settle a precise functional version of that question.
A product of single-shift logarithmic multipliers belongs to the finite
single-shift integer-order jet span exactly in the rational case or when
only one shift is active. The obstruction is analytic monodromy, not a
failed numerical relation search and not a statement about individual
period values.

The positive counterpart is a normally convergent, all-order expansion
of every finite mixed-twist convolution into one central twist. Its
coefficients are finite Pochhammer derivatives, hence generalized harmonic
and Stirling polynomials. Finite Fourier corrections remove the restriction
that the twists lie in one frequency interval. A symmetric specialization
gives explicit Stieltjes contact terms and ordinary convolutions of
rational-grid log-Gamma functions. A unilateral version yields exact
multishift Lerch identities, all ordinary primitives, a meromorphic zeta
continuation, and an all-index formula for logarithmic cutoff constants in
single Hurwitz-zeta derivatives. These constants satisfy a lower-degree
divided-difference recursion and a Gamma antiderivative identity. At the
zero-order spectral intersection, we compute every ray jet: the value is
direction-dependent, whereas the first derivative is a finite additive
log-Gamma expression.

Full proofs, 824 finite exact assertions, and independent numerical
checks accompany the formulas. Classical mechanisms are credited; no
claim of global priority, proof-assistant formalization, or resolution of
fixed-value $S_6/S_8$ identities is made.
\end{abstract}
\vspace{3pt}
\noindent\textbf{Source anchor.} Targeted source reading is pinned to commit\\
\pin.\\
The complete repository and the binary incoming archives were not
exhaustively audited. The package is additive and makes no repository changes.
\clearpage
\tableofcontents
\clearpage

\section{The question and the exact scope of the answer}
\label{sec:scope}
The canonical unequal-twist chapter defines $K_\alpha^\theta$ by the
Fourier multiplier $(2\pi i(n+\theta))^{-\alpha}$, proves a Dirichlet beta
mixture for arbitrary orders in a common frequency interval, and proves
finite partial fractions at positive integer orders \cite{repo-mixed}.
The preceding twisted chapter fixes the meromorphic Lerch family,
its endpoint distributions, its Gamma normalization, and its rational
Hurwitz coordinates \cite{repo-twisted}. We retain those conventions.
The editorial ledger explicitly separates these established results from
finite reduction of mixed Stieltjes jets \cite{repo-ledger}.

There are three distinct questions. A \emph{functional identity} in a
finite span of single-twist integer-order jets is a statement about all
Fourier coefficients. A \emph{convergent single-twist expansion} can use
unbounded positive orders. A \emph{special-value identity} is an equality
at a particular argument and may hold without any functional reduction.
Our nonclosure theorem concerns the first question; our expansion answers
the second. It neither proves nor disproves the third.

The central contributions are Theorem~\ref{thm:central}, its all-jet form
\eqref{eq:jetcompiler}, the finite-closure classification
Theorem~\ref{thm:classification}, the ordinary Gamma convolution
\eqref{eq:gammaconv}, the cutoff formula \eqref{eq:Cformula}, and the
spectral-ray formula \eqref{eq:rayjets}. They form one chain rather than
unrelated examples: all arise by resolving several shifted denominators
about one centre, keeping the logarithmic decorations before taking limits.

The binomial theorem, partial fractions, Fourier uniqueness, the Hurwitz
shift calculus, Gamma products, and cyclotomic filters are classical.
DLMF supplies the relevant normalizations \cite{dlmf-zeta,dlmf-lerch,dlmf-gamma,dlmf-polylog,dlmf-polygamma}.
The broader analytic and multivalued setting of Lerch functions is treated
by Lagarias and Li \cite{lagarias-li}. We do not claim that a literature
priority search is complete. In particular, a derived special case of a
classical Gamma product is identified as such, not advertised as a new
special-function discovery.

\section{Conventions and the normalized twisted family}
\label{sec:conventions}
For $0<\theta<1$ and $n\in\Z$, set
\begin{equation}
 \lambda_{n,\theta}=2\pi i(n+\theta),\qquad
 L_{n,\theta}=\log(2\pi|n+\theta|)
       +\frac{i\pi}{2}\sgn(n+\theta).
 \label{eq:lambda}
\end{equation}
All powers of $\lambda$ use this logarithm. Fourier coefficients are
$\wt f(n)=\int_0^1f(x)e^{-2\pi i n x}\dd x$, and circular convolution
multiplies Fourier coefficients. Define
\begin{equation}
 \wt K_\alpha^\theta(n)=e^{-\alpha L_{n,\theta}},\qquad
 J_m^{r,\theta}=(-1)^m
      \left.\partial_\alpha^mK_\alpha^\theta\right|_{\alpha=r}.
 \label{eq:KJ}
\end{equation}
Thus $\wt J_m^{r,\theta}(n)=L_{n,\theta}^m\lambda_{n,\theta}^{-r}$.
The order $r$ may be complex until an integer assumption is stated.
In particular, $J_0^{0,\theta}=K_0^\theta=\delta_0$; the zero Fourier
coefficient is retained, and is not a zero \emph{shifted} frequency.

On a compact set of orders the multipliers and any fixed number of
order derivatives grow at most polynomially in $|n|$. Pairing with a
smooth periodic test function therefore proves that $K_\alpha^\theta$
is entire as a distribution-valued function. Products of these polynomially
bounded sequences again define distributions, so all convolutions used
below are well-defined. This is not pointwise multiplication of singular
functions.

Write $W_s^\theta=K_{1-s}^\theta$ and
$Z_s^\theta=\Gamma(1-s)W_s^\theta$. The canonical normalization is
\begin{equation}
 Z_{1+u}^\theta=-\frac{\delta_0}{u}
       +\sum_{m\ge0}\frac{(-1)^m}{m!}U_m^\theta u^m.
 \label{eq:Udefinition}
\end{equation}
Away from $x=0$, the family has the ordinary representative
\begin{equation}
 Z_s^\theta(x)=e^{-2\pi i\theta x}
     \Phi(e^{-2\pi i\theta},s,x),\qquad 0<x<1,
 \label{eq:Zrep}
\end{equation}
with the endpoint continuation fixed by \eqref{eq:Udefinition}; see
\cite{repo-twisted}. Expanding $\Gamma(-u)$ gives the useful dictionary
\begin{align}
 E_m^\theta&:=J_m^{0,\theta},\notag\\
 E_1^\theta&=-U_0^\theta-\gamma\delta_0,\label{eq:Edict1}\\
 E_2^\theta&=2U_1^\theta+2\gamma U_0^\theta
                +(\gamma^2-\zeta(2))\delta_0.\label{eq:Edict2}
\end{align}
Here $\gamma$ is Euler's constant. In particular the Gamma normalization
cannot be suppressed without changing point-supported terms.

For scalar functions we use
\begin{equation}
 \zeta(1+u,a)=\frac1u+
       \sum_{m\ge0}\frac{(-1)^m}{m!}\gamma_m(a)u^m,
 \qquad \gamma_0(a)=-\psi(a),\quad a>0.
 \label{eq:stieltjes}
\end{equation}
A superscript $(j)$ on $\zeta$ or $\Phi$ denotes differentiation in the
\emph{order}, not in the final parameter. Harmonic numbers are
$H_N^{(j)}=\sum_{\ell=1}^N\ell^{-j}$, with $H_0^{(j)}=0$ and
$H_N=H_N^{(1)}$. The unsigned Stirling number of the first kind is
$\stir{k}{q}$, with $\stir00=1$ and value zero outside $0\le q\le k$.

\section{A convergent expansion for every finite set of twists}
\label{sec:central}
Let $\theta_1,\ldots,\theta_d\in(0,1)$, choose $c\in(0,1)$, and put
\begin{equation}
 d_j=\theta_j-c,\quad R=\max_j|d_j|,\quad
 \eta=\min(c,1-c),\quad A=\sum_{j=1}^d\alpha_j.
 \label{eq:centraldata}
\end{equation}
Assume $R<\eta$. Such a centre always exists: the midpoint of the least
and greatest twist works. Define the finite polynomials
\begin{align}
 h_k(\av;\dv)
 &=\sum_{k_1+\cdots+k_d=k}\prod_{j=1}^d
       \frac{(\alpha_j)_{k_j}}{k_j!}d_j^{k_j},\label{eq:hk}\\
 \sum_{k\ge0}h_k t^k
 &=\prod_{j=1}^d(1-d_jt)^{-\alpha_j},\label{eq:hgen}\\
 h_0&=1,\qquad
 k h_k=\sum_{\ell=1}^k
       \left(\sum_j\alpha_jd_j^\ell\right)h_{k-\ell}.
 \label{eq:hrec}
\end{align}
The recurrence follows by logarithmic differentiation of
\eqref{eq:hgen}, and supplies an exact coefficient algorithm.

\begin{theorem}[Entire multicentre expansion]\label{thm:central}
For all $\av\in\C^d$,
\begin{equation}
 K_{\alpha_1}^{\theta_1}*\cdots*K_{\alpha_d}^{\theta_d}
   =\sum_{k=0}^\infty(-2\pi i)^k
           h_k(\av;\dv)K_{A+k}^{c}.
 \label{eq:central}
\end{equation}
The series converges locally uniformly in the distribution topology,
may be differentiated any finite number of times in the orders, and
retains every Fourier coefficient, including $n=0$.
\end{theorem}
\begin{proof}
At a fixed mode all $n+\theta_j$ have the same sign as $n+c$. Since
$|d_j/(n+c)|<1$, the real number $1+d_j/(n+c)$ is positive. Hence
\[
 L_{n,\theta_j}=L_{n,c}+\log\left(1+\frac{d_j}{n+c}\right)
\]
with no missing phase. The ordinary binomial theorem gives
\begin{equation}
 \prod_j\lambda_{n,\theta_j}^{-\alpha_j}
 =\lambda_{n,c}^{-A}
     \sum_{k\ge0}\frac{(-1)^kh_k(\av;\dv)}{(n+c)^k}.
 \label{eq:modeexp}
\end{equation}
This is precisely the Fourier coefficient of \eqref{eq:central}.

For normal convergence, choose $R_0$ strictly between $R$ and $\eta$.
The generating function \eqref{eq:hgen} is bounded on $|t|=R_0^{-1}$
uniformly on a compact set of $\av$; Cauchy's estimate gives
$|h_k|\le C R_0^k$. The sum of the absolute values of the terms in
\eqref{eq:modeexp} is therefore bounded by a constant times
$|\lambda_{n,c}^{-A}|$, uniformly in $n$. This last factor has at most
polynomial growth uniformly on the compact order set. Multiply by the
rapidly decreasing Fourier coefficients of a test function and sum.
This proves absolute double-sum convergence and normal convergence of
all pairings. Cauchy's formula in the order variables justifies all
order derivatives. Fourier uniqueness proves the identity.
\end{proof}

\begin{remark}[What the expansion accomplishes]
The series is geometric in the shift displacement, not merely asymptotic.
Its coefficients are finite polynomials, while its single-twist orders
increase without bound. The result is compatible with finite partial
fractions at integer orders, but it is not obtained by differentiating
an identity stated only on the integer lattice.
\end{remark}

\subsection{Arbitrary real nonintegral twists and finite mode corrections}
The common-interval hypothesis can be removed from the expansion, though
not from the unsubtracted beta mixture used in the manuscript.
Let now $\theta_j\in\mathbb R\setminus\Z$ be arbitrary and choose any
$c\in\mathbb R\setminus\Z$. Take an integer $N$ large enough that
\begin{equation}
 |n+c|>R:=\max_j|\theta_j-c|\qquad (|n|\ge N).
 \label{eq:Ncondition}
\end{equation}
For a distribution $F$, let $P_NF$ retain the modes $|n|<N$, and put
$K_\alpha^{c,[N]}=(I-P_N)K_\alpha^c$.

\begin{corollary}[Exact finite-mode repair]\label{cor:alltwists}
With $h_k$ as in \eqref{eq:hk},
\begin{align}
 K_{\alpha_1}^{\theta_1}*\cdots*K_{\alpha_d}^{\theta_d}
 ={}&\sum_{|n|<N}e^{2\pi inx}
            \prod_j\lambda_{n,\theta_j}^{-\alpha_j}\notag\\
 &+\sum_{k\ge0}(-2\pi i)^kh_k(\av;\dv)
                            K_{A+k}^{c,[N]}.
 \label{eq:headrepair}
\end{align}
The second series is normally convergent in distributions and can be
differentiated in every order parameter.
\end{corollary}
\begin{proof}
Condition \eqref{eq:Ncondition} guarantees both the common sign and the
convergent binomial expansion at all retained tail modes. Their denominators
are bounded below by a number strictly greater than $R$. The proof of
Theorem~\ref{thm:central} applies there. The exceptional modes are restored
with their actual logarithms, so no incorrect common phase is assigned
to them.
\end{proof}

\section{All mixed order jets and their harmonic coefficients}
\label{sec:jets}
For multi-indices, write $|\mv|=\sum_jm_j$,
$\binom{\mv}{\qv}=\prod_j\binom{m_j}{q_j}$, and
$\partial^{\qv}=\prod_j\partial_{\alpha_j}^{q_j}$.
Put $R_*:=\sum_jr_j$ and $M:=|\mv|$; $R_*$ is an order sum, not the
geometric radius $R$ of \eqref{eq:centraldata}.

\begin{theorem}[The all-jet coefficient formula]\label{thm:jets}
For nonnegative integers $m_j$ and arbitrary complex $r_j$,
\begin{align}
 J_{m_1}^{r_1,\theta_1}*\cdots*J_{m_d}^{r_d,\theta_d}
 =\sum_{k\ge0}(-2\pi i)^k
   \sum_{\qv\le\mv}(-1)^{|\qv|}\binom{\mv}{\qv}
   (\partial^{\qv}h_k)(\rv;\dv)
     J_{M-|\qv|}^{R_*+k,c}.
 \label{eq:jetcompiler}
\end{align}
The same formula applies to \eqref{eq:headrepair}, with the differentiated
finite head and the corresponding projected central jets.
\end{theorem}
\begin{proof}
Apply $(-1)^M\partial^{\mv}$ to \eqref{eq:central}. In each Leibniz
term, $|\qv|$ derivatives act on $h_k$ and the remaining $M-|\qv|$
derivatives act on $K_{A+k}^c$. The latter give
$(-1)^{M-|\qv|}J_{M-|\qv|}^{R_*+k,c}$; the resulting sign is
$(-1)^{|\qv|}$. The normal convergence already proved justifies the
operation, including at zero and negative integer orders.
\end{proof}

Every coefficient in \eqref{eq:jetcompiler} is a finite sum of products
of
\begin{equation}
 D_q(r,k)=\left.\partial_u^q\frac{(u)_k}{k!}\right|_{u=r}.
 \label{eq:Ddef}
\end{equation}
For positive integer $r$, let
$\mathcal H_j=H_{r+k-1}^{(j)}-H_{r-1}^{(j)}$.
The exponential complete Bell polynomial $B_q$ is specified by
$\exp(\sum_{j\ge1}x_jt^j/j!)=\sum_{q\ge0}B_q(x_1,\ldots,x_q)t^q/q!$.
Then
\begin{equation}
 D_q(r,k)=\frac{(r)_k}{k!}
 B_q\bigl(\mathcal H_1,-1!\mathcal H_2,2!\mathcal H_3,
                 \ldots,(-1)^{q-1}(q-1)!\mathcal H_q\bigr).
 \label{eq:Dpositive}
\end{equation}
This follows by taking logarithms of
$(r+u)_k/(r)_k=\prod_{j=0}^{k-1}(1+u/(r+j))$.
At the important zero order one must not divide by $(0)_k$:
\begin{equation}
 D_q(0,k)=\frac{q!}{k!}\stir{k}{q},\qquad k,q\ge0.
 \label{eq:Dzero}
\end{equation}
In particular, for $k\ge1$,
\begin{align}
 D_1(0,k)&=\frac1k,&
 D_2(0,k)&=\frac{2H_{k-1}}k,\notag\\
 D_3(0,k)&=\frac{3(H_{k-1}^2-H_{k-1}^{(2)})}{k}.
 \label{eq:Dlow}
\end{align}
These identities also follow from
$(u)_k/k!=u\prod_{j=1}^{k-1}(1+u/j)/k$.
They explain why differentiating before setting an order to zero is
essential: the value $(0)_k=0$ does not imply that its jets vanish.

\subsection{A symmetric first-jet identity with its contact term}
Set $a=c-d$, $b=c+d$, with $|d|<\min(c,1-c)$, and define
\begin{equation}
 \eta_k=H_{2k-1}-H_{k-1},\qquad k\ge1.
 \label{eq:etak}
\end{equation}

\begin{corollary}[Symmetric mixed Stieltjes generators]\label{cor:symmetric}
The following identity holds in periodic distributions:
\begin{equation}
 E_1^a*E_1^b=E_2^c-
  \sum_{k\ge1}\frac{(2\pi i d)^{2k}}k
       \left(J_1^{2k,c}+\eta_kK_{2k}^c\right).
 \label{eq:Esymmetric}
\end{equation}
In the canonical unnormalized Stieltjes coordinates this is
\begin{align}
 U_0^a*U_0^b={}&2U_1^c+2\gamma U_0^c
          -\gamma(U_0^a+U_0^b)-\zeta(2)\delta_0\notag\\
 &-\sum_{k\ge1}\frac{(2\pi i d)^{2k}}k
       \left(J_1^{2k,c}+\eta_kK_{2k}^c\right).
 \label{eq:Usymmetric}
\end{align}
\end{corollary}
\begin{proof}
At a mode put $L=L_{n,c}$ and $u=d/(n+c)$. Then
\[
 L_{n,a}L_{n,b}=L^2+L\log(1-u^2)
                       +\log(1-u)\log(1+u).
\]
The elementary power-series identity
\begin{equation}
 \log(1-u)\log(1+u)
       =-\sum_{k\ge1}\frac{\eta_k}{k}u^{2k}
 \label{eq:logproduct}
\end{equation}
follows from
$\log^2(1-u)=\sum_{j\ge2}2H_{j-1}u^j/j$ and
$2\log(1-u)\log(1+u)=\log^2(1-u^2)-\log^2(1-u)-\log^2(1+u)$.
This proves \eqref{eq:Esymmetric} mode by mode. Substituting
\eqref{eq:Edict1}--\eqref{eq:Edict2} gives \eqref{eq:Usymmetric}.
The first correction term in \eqref{eq:Esymmetric} is
$-(2\pi id)^2(J_1^{2,c}+K_2^c)$; this checks both signs.
\end{proof}

When $d=0$, \eqref{eq:Usymmetric} becomes
$U_0^c*U_0^c=2U_1^c-\zeta(2)\delta_0$, in agreement with the
one-twist Gamma normalization. Omitting the point mass gives an
incorrect identity even at the zero Fourier coefficient.

\section{A complete finite-closure classification in the integer ladder}
\label{sec:obstruction}
The next theorem concerns a \emph{specified function space}. Coefficients
in a linear combination are constants with respect to the frequency
variable; they may depend on the fixed shifts. A pointwise evaluation
of the resulting functions is a different problem.

\begin{definition}
Let $\mathcal V$ consist of all finite linear combinations of sequences
\begin{equation}
 (n+c)^{-j}\log^q(n+c),\qquad j\in\Z,\quad q\in\Z_{\ge0},
 \label{eq:span}
\end{equation}
where finitely many centres $c$ are allowed and all expressions are
understood for sufficiently large positive integers $n$ on their
standard large-positive branch. Equality is allowed to fail at finitely
many $n$. The centres need not be prescribed in advance.
\end{definition}
The allowance $j\le0$ includes polynomials and hence the Fourier
multipliers of point masses and their derivatives at the origin.
The factors $2\pi i$ and the additive logarithm constant in
\eqref{eq:lambda} do not change this span: a finite binomial change of
logarithmic coordinates absorbs them.

\begin{lemma}[From eventual sequence equality to a function identity]
\label{lem:sequence}
Suppose $F(y)$ is a finite sum of products of integer powers of $y+c$
and nonnegative integer powers of $\log(y+c)$. If $F(n)=0$ for every
sufficiently large positive integer, its analytic germ at large positive
$y$ is identically zero.
\end{lemma}
\begin{proof}
For sufficiently large $|y|$ in a fixed sector about the positive real
axis, expand each shifted power and write
$\log(y+c)=\log y+\log(1+c/y)$. Thus
\[
 F(y)=\sum_{p\ge p_0}y^{-p}P_p(\log y),
\]
where the coefficients of each power of $\log y$ form convergent
Laurent series in $1/y$, and the degrees of all $P_p$ have a common
finite bound $D$. If not all $P_p$ vanish, choose the least $p$ with
$P_p\ne0$. At positive integers,
\[
 F(n)=n^{-p}\{P_p(\log n)+O(n^{-1}(\log n)^D)\}.
\]
A nonzero polynomial in $\log n$ dominates the remainder for all
sufficiently large $n$. This contradicts eventual vanishing. Therefore
all Laurent coefficients vanish and the analytic germ is zero.
\end{proof}

\begin{theorem}[Finite single-centre closure criterion]
\label{thm:classification}
Let $a_1,\ldots,a_d$ be distinct real shifts, and let $r_j,m_j$ be
nonnegative integers. Remove inactive sites with $r_j=m_j=0$.
Then the sequence
\begin{equation}
 F(n)=\prod_{j=1}^d(n+a_j)^{-r_j}\log^{m_j}(n+a_j)
 \label{eq:target}
\end{equation}
belongs to $\mathcal V$ if and only if either all $m_j=0$, or there
is at most one active site.
\end{theorem}
\begin{proof}
If all $m_j=0$, ordinary partial fractions give the required finite
representation. With one active site the target is itself a generator.

Conversely, an eventual sequence representation gives an analytic
function identity by Lemma~\ref{lem:sequence}. Continue this identity
along loops around the finitely many points $-a_j$ and the candidate
centres. Let $\Delta_a$ denote continuation once counterclockwise
around $-a$ minus the original germ. It replaces $\log(y+a)$ by
$\log(y+a)+2\pi i$, and fixes logarithms at every other centre and
all integer rational powers. These operations commute on the functions
under consideration.

Suppose two sites, say $a$ and $b$, have positive logarithmic degree.
Then $\Delta_a\Delta_bF$ is nonzero: it is the nonzero rational
prefactor times the product of the two nonzero finite differences of
positive-degree logarithm powers and the other factors. Every generator
in \eqref{eq:span}, however, is killed by $\Delta_a\Delta_b$. This
contradicts the continued identity.

It remains to consider exactly one logarithmically decorated site $a$,
of degree $m\ge1$, and another active site $b$ with $r_b>0$.
Write $F(y)=R(y)\log^m(y+a)$, where $R$ is rational and has a pole
at $-b$. Group the putative single-centre representation by centre and
log degree. Its $a$-centred portion is
$\sum_qR_q(y)\log^q(y+a)$, with each $R_q$ a Laurent polynomial in
$y+a$. All other portions have trivial monodromy around $-a$.
Because $\Delta_a^{m+1}F=0$, this identity forces $R_q=0$ for $q>m$.
Indeed, a polynomial in $\log(y+a)$ with Laurent-polynomial coefficients
has a unique such expansion, by the same leading-power argument near
$y=-a$ as in Lemma~\ref{lem:sequence}; its highest log degree cannot
cancel a lower one under repeated finite differences.
Applying $\Delta_a^m$ now yields
\[
 (2\pi i)^m m!R(y)=(2\pi i)^m m!R_m(y).
\]
The right side can have no finite pole except at $-a$, whereas the left
has a pole at $-b$. This is impossible.
\end{proof}

\begin{corollary}[Mixed twisted jets do not generally close finitely]
\label{cor:nonclosure}
For distinct nonintegral twists, a convolution of jets
$J_{m_j}^{r_j,\theta_j}$ with nonnegative integer orders satisfies the
same criterion for membership in any finite constant-coefficient span
of single-twist integer-order jets. Finite Fourier corrections and
point masses or their derivatives at the origin cannot repair the
obstruction. In particular $E_1^a*E_1^b$, and also $U_0^a*U_0^b$, do
not admit such a finite reduction when $a\ne b$.
\end{corollary}
\begin{proof}
Apply the theorem to positive Fourier modes. Scaling the powers by
$2\pi i$ and translating the logarithms by $\log(2\pi)+i\pi/2$ are
finite triangular transformations. For the $U_0$ product,
\eqref{eq:Edict1} shows that its difference from $E_1^a*E_1^b$ lies
in the single-twist span. Equality modulo finite Fourier support is
already covered by eventual sequence equality.
\end{proof}

\paragraph{Limits of the conclusion.}
This is a functional nonclosure result for the integer-order ladder,
not a proof of numerical period independence, not non-elementarity at
a fixed $x$, and not a theorem excluding arbitrary argument-dependent
coefficients, arbitrary complex-order bases, or every possible
special-function enlargement. In particular, evaluating at a rational
$x$ can introduce cancellations which the theorem does not address.
The infinite expansion \eqref{eq:jetcompiler} is not in conflict with
nonclosure: it uses infinitely many integer orders.

\section{Covariant primitives and ordinary log-Gamma convolutions}
\label{sec:gamma}
For $D_c=\partial_x+2\pi ic$, the nonintegrality of $c$ makes $D_c$
invertible on periodic distributions. Its inverse of order $p\ge0$
has multiplier $\lambda_{n,c}^{-p}$ and hence
\begin{equation}
 D_c^{-p}J_m^{r,c}=J_m^{r+p,c}.
 \label{eq:primitive}
\end{equation}
It may be applied termwise to \eqref{eq:jetcompiler} or
\eqref{eq:Esymmetric}. These are \emph{covariant} primitives, not
an unqualified inverse of the ordinary periodic derivative. If
$D_cu=f$, then $(e^{2\pi icx}u)'=e^{2\pi icx}f$, so the integrating
factor converts them into ordinary weighted antiderivatives.

For positive integer $r$, the ordinary representative of a central
kernel is
\begin{equation}
 K_r^c(x)=\frac{e^{-2\pi icx}}{\Gamma(r)}
               \Phi(e^{-2\pi ic},1-r,x),\qquad 0<x<1.
 \label{eq:Krrep}
\end{equation}
Differentiating the entire order identity gives
\begin{align}
 J_1^{r,c}(x)&=\frac{e^{-2\pi icx}}{\Gamma(r)}
       \left(\Phi^{(1)}+\psi(r)\Phi\right),\label{eq:J1rep}\\
 J_2^{r,c}(x)&=\frac{e^{-2\pi icx}}{\Gamma(r)}
       \left(\Phi^{(2)}+2\psi(r)\Phi^{(1)}
                    +(\psi(r)^2-\psi'(r))\Phi\right),
 \label{eq:J2rep}
\end{align}
where every $\Phi$ on the right is evaluated at
$(e^{-2\pi ic},1-r,x)$. At integer $r$, the coefficients reduce to
$\psi(r)=H_{r-1}-\gamma$ and
$\psi'(r)=\zeta(2)-H_{r-1}^{(2)}$. All higher jets follow by the same
Leibniz rule for $1/\Gamma(r)$; the general master derivative avoids
any division by $\Gamma(0)$ at the exceptional order zero.

\subsection{Finite rational-grid Gamma representatives}
Let $\theta=p/q\in(0,1)$ be rational and $z=e^{-2\pi ip/q}$. Splitting
an ordinary Lerch series into residues gives
\begin{equation}
 \Phi(z,s,x)=q^{-s}\sum_{\ell=0}^{q-1}
            z^\ell\zeta\left(s,\frac{x+\ell}{q}\right).
 \label{eq:Lerchrational}
\end{equation}
This extends to all $s$; the poles of the individual Hurwitz functions
cancel because $\sum z^\ell=0$. Using
$\zeta'(0,y)=\log\Gamma(y)-\tfrac12\log(2\pi)$ and
$\Phi(z,0,x)=(1-z)^{-1}$ gives an explicit representative for
$P_\theta:=J_1^{1,\theta}$:
\begin{equation}
 P_{p/q}(x)=e^{-2\pi ipx/q}\left\{
   \sum_{\ell=0}^{q-1}z^\ell\log\Gamma\left(\frac{x+\ell}{q}\right)
       -\frac{\gamma+\log q}{1-z}\right\}.
 \label{eq:Pgamma}
\end{equation}
The real positive branch of $\log\Gamma$ is used on these arguments.
The $\log q$ and Euler-constant term are necessary, not optional
normalization choices.

The Fourier coefficients of $P_\theta$ are $L_{n,\theta}/\lambda_{n,\theta}$,
which are square summable. Thus $P_\theta\in L^2(\T)$, and the circular
convolution of any two such functions is an ordinary continuous
function. The integral is not a finite part.

\begin{theorem}[An ordinary mixed Gamma convolution]\label{thm:gammaconv}
Let $a=c-d$, $b=c+d$, with $|d|<\min(c,1-c)$. Set
\begin{equation}
 S_0=0,\qquad S_k=\sum_{\ell=1}^k\frac{\eta_\ell}{\ell}.
 \label{eq:Sk}
\end{equation}
Then, for every $x\in\T$,
\begin{align}
 \int_0^1P_a(t)P_b(\{x-t\})\dd t
   =\sum_{k=0}^\infty(2\pi id)^{2k}
       \left(J_2^{2k+2,c}(x)-H_kJ_1^{2k+2,c}(x)
                          -S_kK_{2k+2}^c(x)\right).
 \label{eq:gammaconv}
\end{align}
The series is uniformly convergent. If $a,b$ are rational,
\eqref{eq:Pgamma} makes its left side an explicit integral of finite
log-Gamma combinations. If $c$ is rational,
\eqref{eq:Lerchrational} makes every term on the right a finite
Hurwitz order-jet expression at a negative integer.
\end{theorem}
\begin{proof}
At a Fourier mode, divide the logarithm product in
\eqref{eq:logproduct} by
$\lambda_{n,a}\lambda_{n,b}=\lambda_{n,c}^2(1-u^2)$.
Multiplication by $(1-u^2)^{-1}$ yields the coefficient of $u^{2k}$
\[
 L_{n,c}^2-H_kL_{n,c}-S_k.
\]
This proves the Fourier coefficient identity in \eqref{eq:gammaconv}.
The sum of the absolute coefficient contributions, after summing over
$k$, is bounded by a constant times
$(1+\log^2(2+|n|))/(1+|n|)^2$. Summing over $n$ proves absolute
double convergence and uniform convergence in $x$. The $L^2$
convolution has exactly these Fourier coefficients and is continuous;
Fourier uniqueness proves equality everywhere.
\end{proof}

\subsection{Two half-grid zeta-series evaluations}
For a particularly explicit family put $c=\tfrac12$, $|d|<\tfrac12$,
$a=\tfrac12-d$, $b=\tfrac12+d$, and $\ell=\log(2\pi)$. Write
\[
 Z(s)=\zeta(s,\tfrac12)=(2^s-1)\zeta(s),\qquad
 A(s)=\Phi(-1,s,\tfrac12)
       =2^{-s}\bigl(\zeta(s,\tfrac14)-\zeta(s,\tfrac34)\bigr).
\]
Theorem~\ref{thm:gammaconv} specializes to the ordinary integral identities
\begin{align}
 \int_0^1P_a(t)P_b(1-t)\dd t
 ={}&-\frac1{2\pi^2}\sum_{k\ge0}d^{2k}\Bigl[Z''(2k+2)
          +(H_k-2\ell)Z'(2k+2)\notag\\
 &\hspace{17mm}+\bigl(\ell^2-H_k\ell-S_k-\tfrac{\pi^2}4\bigr)Z(2k+2)\Bigr],
 \label{eq:gammazero}\\
 \int_0^1P_a(t)P_b(\{\tfrac12-t\})\dd t
 ={}&\frac{i}{4\pi}\sum_{k\ge0}d^{2k}
       \Bigl[2A'(2k+2)\notag\\
 &\hspace{27mm}+(H_k-2\ell)A(2k+2)\Bigr].
 \label{eq:gammahalf}
\end{align}
For rational $a,b$, these are explicit log-Gamma integrals in positive
integer zeta derivatives. They are new derived instances in this
continuation, not assertions of first occurrence in the literature.

To verify the formulas, pair the modes $n$ and $-n-1$:
\begin{align*}
 K_s^{1/2}(0)&=2\cos(\pi s/2)(2\pi)^{-s}Z(s),\\
 K_s^{1/2}(\tfrac12)&=-2i\sin(\pi s/2)(2\pi)^{-s}A(s).
\end{align*}
Differentiate these identities in $s$ once with a minus sign and twice
with a plus sign, and evaluate at $s=2k+2$. The second derivative of
the cosine contributes the $-\pi^2Z/4$ term in \eqref{eq:gammazero}.
This term would be lost by specializing the phase before differentiating.

\section{The unilateral multishift Lerch family and its primitives}
\label{sec:unilateral}
For positive shifts $a_1,\ldots,a_d$, define initially
\begin{equation}
 \mathcal L(z;\av,\boldsymbol a)
   =\sum_{n=0}^\infty z^n\prod_j(n+a_j)^{-\alpha_j}.
 \label{eq:multilerch}
\end{equation}
For $|z|<1$ it is entire in the orders. At $z=1$ the defining series
converges absolutely when $\Re A>1$.
Choose $c>0$ and $d_j=a_j-c$ with $R=\max|d_j|<c$; the midpoint of
the extreme positive shifts again works.

\begin{theorem}[Single-Lerch expansion and meromorphic zeta continuation]
\label{thm:multizeta}
For $|z|<1$,
\begin{equation}
 \mathcal L(z;\av,\boldsymbol a)
    =\sum_{k\ge0}(-1)^kh_k(\av;\dv)\Phi(z,A+k,c).
 \label{eq:unilateral}
\end{equation}
This is normally convergent on compact order sets and compact subsets
of the disk. At $z=1$ it gives the meromorphic continuation
\begin{equation}
 \mathcal Z(\av;\boldsymbol a)
    =\sum_{k\ge0}(-1)^kh_k(\av;\dv)\zeta(A+k,c).
 \label{eq:multizeta}
\end{equation}
The only possible polar hyperplanes are $A=1-k$, $k\ge0$, with
residue in the coordinate $A-(1-k)$ equal to
\begin{equation}
 \left.(-1)^kh_k(\av;\dv)\right|_{A=1-k}.
 \label{eq:residues}
\end{equation}
A vanishing residue removes the corresponding pole; the theorem does
not assert that every possible hyperplane is an actual divisor.
\end{theorem}
\begin{proof}
Expand the product in \eqref{eq:multilerch} by
\eqref{eq:hgen} with $t=-1/(n+c)$. In the disk, exponential damping
justifies the double sum locally uniformly. At $z=1$, interchange is
first justified for $\Re A>1$. For a compact order set and sufficiently
large $k$, $|\zeta(A+k,c)|\le Cc^{-k}$; finitely many earlier terms
are treated individually. The bound $|h_k|\le CR_0^k$ for a fixed
$R<R_0<c$ proves normal convergence away from the finitely many
hyperplanes meeting the compact set. Only the $k$th term can have a
pole at $A=1-k$, and the Hurwitz residue is one. This proves the stated
continuation and residue formula.
\end{proof}

Alternatively one may use any positive centre after taking a finite
head. Explicitly, for $N$ with $\max_j|a_j-c|<N+c$,
\begin{align}
 \mathcal L(z;\av,\boldsymbol a)
 ={}&\sum_{n=0}^{N-1}z^n\prod_j(n+a_j)^{-\alpha_j}\notag\\
 &+z^N\sum_{k\ge0}(-1)^kh_k(\av;\boldsymbol a-c)
                         \Phi(z,A+k,N+c).
 \label{eq:unilateralhead}
\end{align}
This gives the same continuation and is useful when a preferred
arithmetic centre is fixed.

\subsection{Differential relations and all ordinary primitives}
Termwise operations in the disk give the exact relations
\begin{align}
 (z\partial_z+a_j)\mathcal L(z;\av,\boldsymbol a)
       &=\mathcal L(z;\av-\boldsymbol e_j,\boldsymbol a),\label{eq:eulerderivative}\\
 \partial_{a_j}\mathcal L(z;\av,\boldsymbol a)
       &=-\alpha_j\mathcal L(z;\av+\boldsymbol e_j,\boldsymbol a).
 \label{eq:shiftderivative}
\end{align}
Consequently, on a chosen branch of $z^{a_j}$,
\begin{equation}
 \int_0^z w^{a_j-1}\mathcal L(w;\av,\boldsymbol a)\dd w
    =z^{a_j}\mathcal L(z;\av+\boldsymbol e_j,\boldsymbol a).
 \label{eq:eulerprimitive}
\end{equation}
There is also an unweighted formula with no change of variable.
For each positive integer $p$,
\begin{align}
 \frac1{(p-1)!}\int_0^z(z-w)^{p-1}
                \mathcal L(w;\av,\boldsymbol a)\dd w
  =z^p\mathcal L\bigl(z; (\av,\underbrace{1,\ldots,1}_{p});
                     (\boldsymbol a,1,2,\ldots,p)\bigr).
 \label{eq:ordinaryprimitives}
\end{align}
Indeed, integrating $w^n$ produces
$z^{n+p}n!/(n+p)!=z^{n+p}\prod_{j=1}^p(n+j)^{-1}$.
All order derivatives commute with these operations on compact subsets
of the disk. Thus \eqref{eq:ordinaryprimitives} supplies exact ordinary
antiderivatives of every logarithmically decorated multishift series.
Any repeated shift can be combined by adding its orders.

\section{Logarithmic cutoff constants: all indices and exact recursions}
\label{sec:cutoff}
For $a,b>0$ and nonnegative integers $m,n$, let $M=m+n$ and define
\begin{equation}
 C_{m,n}(a,b)=\lim_{N\to\infty}\left\{
  \sum_{j=0}^{N-1}\frac{\log^m(j+a)\log^n(j+b)}{j+a}
          -\frac{\log^{M+1}N}{M+1}\right\}.
 \label{eq:Cdef}
\end{equation}
This is a sharp-cutoff finite part. No Abel or spectral-ray convention
is silently substituted.

\begin{theorem}[A convergent all-index Hurwitz-jet formula]
\label{thm:cutoff}
The limit \eqref{eq:Cdef} exists and is smooth in positive $a,b$.
If $|b-a|<a$, then
\begin{align}
 C_{m,n}(a,b)=\gamma_M(a)+(-1)^M
   \sum_{k=1}^\infty (a-b)^k
    \sum_{q=1}^{\min(n,k)}\binom nq
       \frac{q!}{k!}\stir{k}{q}\,
                \zeta^{(M-q)}(k+1,a).
 \label{eq:Cformula}
\end{align}
The series is locally normally convergent in the stated shift domain.
For general positive $a,b$, choose an integer $N$ with
$|b-a|<a+N$ and use the exact finite-head identity
\begin{equation}
 C_{m,n}(a,b)=\sum_{j=0}^{N-1}
          \frac{\log^m(j+a)\log^n(j+b)}{j+a}
                       +C_{m,n}(a+N,b+N).
 \label{eq:Chead}
\end{equation}
In particular, $C_{m,0}(a,b)=\gamma_m(a)$ and
$C_{m,n}(a,a)=\gamma_M(a)$.
\end{theorem}
\begin{proof}
Subtract the diagonal summand. For $M\ge1$, the difference is
$O(j^{-2}(1+\log^{M-1}j))$, locally uniformly with its fixed-order
parameter derivatives on compact positive shift sets. For $M=0$ it
is zero. The ordinary Stieltjes cutoff limit then proves existence and
smoothness, with the absolutely convergent expression
\begin{equation}
 C_{m,n}(a,b)=\gamma_M(a)+\sum_{j\ge0}
   \frac{\log^m(j+a)\{\log^n(j+b)-\log^n(j+a)\}}{j+a}.
 \label{eq:Cconvergent}
\end{equation}
Replacing $\log N$ by $\log(N+a)$ in the diagonal subtraction changes
the limit by zero.

For the formula, specialize \eqref{eq:multizeta} to two shifts, with
centre $a$ and $d=b-a$:
\begin{equation}
 \mathcal Z(\alpha,\beta;a,b)
   =\zeta(\alpha+\beta,a)+
       \sum_{k\ge1}\frac{(-d)^k(\beta)_k}{k!}
                         \zeta(\alpha+\beta+k,a).
 \label{eq:twozeta}
\end{equation}
Subtract $1/(\alpha+\beta-1)$. The remainder is holomorphic near
$(\alpha,\beta)=(1,0)$. Its
$(-1)^M\partial_\alpha^m\partial_\beta^n$ derivative there equals
\eqref{eq:Cdef}: the uniform leading tail is
$N^{1-\alpha-\beta}/(\alpha+\beta-1)$, and the difference between
this and the subtracted pole gives exactly
$\log^{M+1}N/(M+1)$ after differentiation. The remainder and its jets
tend to zero, by the same $j^{-2}$ subtraction as in
\eqref{eq:Cconvergent}.

The first term contributes $\gamma_M(a)$. For $k\ge1$, the
undifferentiated $(\beta)_k$ vanishes at zero, so at least one of the
$n$ derivatives must hit this factor. Leibniz's rule and
\eqref{eq:Dzero} give \eqref{eq:Cformula}. Normal convergence from
Theorem~\ref{thm:multizeta} justifies the differentiation. Finally,
split the sum defining \eqref{eq:Cdef} at $N$. The difference between
the logarithmic subtractions for the total cutoff and the translated
tail cutoff tends to zero, proving \eqref{eq:Chead}.
\end{proof}

Two useful subfamilies are
\begin{align}
 C_{m,1}(a,b)
 &=\gamma_{m+1}(a)+(-1)^{m+1}
       \sum_{k\ge1}\frac{(a-b)^k}{k}\zeta^{(m)}(k+1,a),
 \label{eq:Cm1}\\
 C_{m,2}(a,b)
 &=\gamma_{m+2}(a)+2(-1)^m\sum_{k\ge1}\frac{(a-b)^k}{k}
       \left(\zeta^{(m+1)}(k+1,a)
                        +H_{k-1}\zeta^{(m)}(k+1,a)\right).
 \label{eq:Cm2}
\end{align}
The second logarithmic decoration therefore introduces harmonic
coefficients even though the zeta values on the right are all single
Hurwitz derivatives.

\begin{theorem}[Lower-degree divided differences and the collision value]
\label{thm:Crec}
For $b\ne a$ and $n\ge1$,
\begin{equation}
 \partial_b C_{m,n}(a,b)
 =\frac{n}{b-a}
        \left(C_{m,n-1}(a,b)-C_{n-1,m}(b,a)\right).
 \label{eq:Crec}
\end{equation}
The apparent collision singularity is removable, with
\begin{equation}
 \left.\partial_b C_{m,n}(a,b)\right|_{b=a}
       =n(-1)^{M-1}\zeta^{(M-1)}(2,a).
 \label{eq:Ccollision}
\end{equation}
Consequently all $C_{m,n}$ are determined recursively in total
logarithmic degree by
\begin{equation}
 C_{m,n}(a,b)=\gamma_M(a)+n\int_a^b
    \frac{C_{m,n-1}(a,t)-C_{n-1,m}(t,a)}{t-a}\dd t.
 \label{eq:Cintegralrec}
\end{equation}
\end{theorem}
\begin{proof}
Differentiate \eqref{eq:Cconvergent}, or the defining cutoff whose
subtraction is independent of $b$. The derivative is the absolutely
convergent sum
\[
 n\sum_{j\ge0}\frac{\log^m(j+a)\log^{n-1}(j+b)}{(j+a)(j+b)}.
\]
Apply $1/(XY)=(1/X-1/Y)/(Y-X)$. The two sums on the right have
identical cutoff divergences, which cancel, yielding \eqref{eq:Crec}.
Setting $b=a$ directly in the convergent differentiated series gives
\eqref{eq:Ccollision}. Integrating from the diagonal gives
\eqref{eq:Cintegralrec}. Both constants on the right have total
logarithmic degree $M-1$, so the recurrence is genuinely lower-degree.
\end{proof}

\subsection{A Gamma antiderivative and polygamma collision jets}
The first decorated constant has the exact formulas
\begin{align}
 C_{0,1}(a,b)&=\gamma_1(a)+\int_a^b
          \frac{\psi(t)-\psi(a)}{t-a}\dd t,\label{eq:C01int}\\
 \int_a^b C_{0,1}(a,t)\dd t
   &=(b-a)C_{0,1}(a,b)-\log\Gamma(b)+\log\Gamma(a)
                                      +(b-a)\psi(a).
 \label{eq:C01primitive}
\end{align}
The first identity is \eqref{eq:Cintegralrec} with $m=0,n=1$;
the second follows by integration by parts, since
$(t-a)\partial_tC_{0,1}(a,t)=\psi(t)-\psi(a)$.
For every $j\ge1$, direct differentiation also gives
\begin{equation}
 \left.\partial_b^j C_{0,1}(a,b)\right|_{b=a}
       =(-1)^{j-1}(j-1)!\zeta(j+1,a)
       =\frac{\psi^{(j)}(a)}{j}.
 \label{eq:C01polygamma}
\end{equation}
Thus this two-shift finite part has a complete polygamma Taylor jet
on the diagonal, while its integral in the second shift has an
explicit log-Gamma correction.

\section{Spectral rays at the zero-order intersection}
\label{sec:rays}
The meromorphic family $\mathcal Z$ is generally not jointly regular
at the origin of its order variables. This is not a defect of the
continuation: the hyperplane $A=0$ can pass through that point.
Nevertheless each transverse ray has a removable one-variable
singularity and a completely explicit jet.

Let $w_1,\ldots,w_d\in\C$, $W=\sum_jw_j\ne0$, and define
\[
 \mathcal Z_{\wv}(t)=\mathcal Z(t\wv;\boldsymbol a),\quad
 D=\sum_jw_j(a_j-c),\quad H_k(t)=h_k(t\wv;\boldsymbol a-c).
\]
The logarithms of $n+a_j$ remain real before multiplication by the
possibly complex weights. This convention avoids an implicit change
of branch for the product.

\begin{theorem}[Every transverse spectral-ray jet]\label{thm:rays}
The singularity of $\mathcal Z_{\wv}$ at $t=0$ is removable, and
\begin{align}
 \mathcal Z_{\wv}(0)&=\frac12-\frac{\sum_jw_ja_j}{W},
 \label{eq:rayvalue}\\
 \mathcal Z_{\wv}'(0)&=\sum_jw_j\log\Gamma(a_j)
                                      -\frac W2\log(2\pi).
 \label{eq:rayGamma}
\end{align}
For every integer $m\ge1$,
\begin{align}
 \mathcal Z_{\wv}^{(m)}(0)
 ={}&W^m\zeta^{(m)}(0,c)
        +(-1)^m mD W^{m-1}\gamma_{m-1}(c)\notag\\
 &+\sum_{k\ge2}(-1)^k
     \sum_{q=1}^{\min(m,k)}\binom mq
             H_k^{(q)}(0)W^{m-q}\zeta^{(m-q)}(k,c).
 \label{eq:rayjets}
\end{align}
The infinite series is normally convergent for the chosen centre.
\end{theorem}
\begin{proof}
The terms $k=0,1$ of \eqref{eq:multizeta} are
\[
 \zeta(Wt,c)-tD\zeta(1+Wt,c).
\]
The second term has the convergent expansion
\begin{equation}
 -\frac DW-D\sum_{j\ge0}\frac{(-1)^j}{j!}
                         \gamma_j(c)W^jt^{j+1}.
 \label{eq:rayexception}
\end{equation}
For $k\ge2$, $\zeta(Wt+k,c)$ is holomorphic near zero and
$H_k(0)=0$. Thus the constant term is
$\zeta(0,c)-D/W=\tfrac12-\sum_jw_ja_j/W$.
Differentiating \eqref{eq:rayexception} and applying Leibniz's rule
in the remaining normally convergent series gives
\eqref{eq:rayjets}.

At first order,
\[
 H_k'(0)=\frac1k\sum_jw_j(a_j-c)^k\quad(k\ge1),
\]
by \eqref{eq:Dzero}. Therefore the first derivative is
\[
 W\zeta'(0,c)-D\gamma_0(c)
      +\sum_jw_j\sum_{k\ge2}\frac{(-1)^k(a_j-c)^k}{k}\zeta(k,c).
\]
The convergent Gamma Taylor formula
\[
 \log\Gamma(c+d)-\log\Gamma(c)-d\psi(c)
       =\sum_{k\ge2}\frac{(-1)^kd^k}{k}\zeta(k,c)
\]
now gives \eqref{eq:rayGamma}, using
$\gamma_0(c)=-\psi(c)$ and
$\zeta'(0,c)=\log\Gamma(c)-\tfrac12\log(2\pi)$.
\end{proof}

The first derivative is independent of the auxiliary centre and is
linear in the ray weights. It yields the zeta-regularized product identity
\begin{equation}
 \exp\bigl(-\mathcal Z_{\wv}'(0)\bigr)
       =\frac{(2\pi)^{W/2}}{\prod_j\Gamma(a_j)^{w_j}}.
 \label{eq:regularizedproduct}
\end{equation}
This is a Gamma-product consequence of the one-dimensional shifted
sequence; it does not assert a general absence of multiplicative
anomalies for other spectra.

In contrast, \eqref{eq:rayvalue} depends on the ray whenever the shifts
are not all equal. For $(a,b)=(1,2)$ the rays $(\alpha,\beta)=(t,t)$
and $(2t,t)$ have limits $-1$ and $-5/6$, respectively. There is
therefore no joint value at $(0,0)$ in this example, even though the
first ray derivative has the additive Gamma formula. The hypothesis
$W\ne0$ is essential: a ray with $W=0$ lies in the possible polar
hyperplane and is not covered by the theorem.

\section{Cyclotomic coordinates and explicit scalar instances}
\label{sec:cyclotomic}
For $1\le r\le q$, let $\omega=e^{2\pi i/q}$ and choose $w$ with
$w^q=z$, $0<|z|<1$. A root-of-unity filter gives
\begin{equation}
 \Phi(z,s,r/q)=q^{s-1}w^{-r}
       \sum_{h=0}^{q-1}\omega^{-rh}\Li_s(\omega^h w).
 \label{eq:rootfilter}
\end{equation}
Indeed the inner sum selects positive integers congruent to $r$ modulo
$q$. Changing the chosen root $w$ permutes the summands and cancels
against $w^{-r}$, so the expression is independent of this choice.
The apparent singularity at $z=0$ is removable. On compact subsets of
the disk all order derivatives may be taken termwise; at derivative
order $m$ this gives
\begin{equation}
 \Phi^{(m)}(z,s,r/q)=q^{s-1}w^{-r}
  \sum_{h=0}^{q-1}\omega^{-rh}
    \sum_{j=0}^m\binom mj(\log q)^{m-j}
                        \Li_s^{(j)}(\omega^hw).
 \label{eq:rootjets}
\end{equation}
The derivative on $\Li_s^{(j)}$ is again with respect to $s$.

At $z=1$, start in $\Re s>1$ and then continue. One must keep the
$h=0$ zeta pole before expanding at $s=1$. With $L=\log q$,
\begin{align}
 \gamma_m(r/q)=(-1)^m\Bigg\{\frac{L^{m+1}}{m+1}
  +\sum_{j=0}^m\binom mj L^{m-j}
   \left[(-1)^j\gamma_j+
       \sum_{h=1}^{q-1}\omega^{-rh}\Li_1^{(j)}(\omega^h)\right]\Bigg\}.
 \label{eq:Stieltjesfilter}
\end{align}
This includes the conductor-log contribution of the pole. It follows
by expanding $q^{s-1}\zeta(s)$, not by treating every summand as
holomorphic at one.

Equations \eqref{eq:unilateral}, \eqref{eq:Cformula}, and
\eqref{eq:rootjets} express rational-shift mixed logarithmic sums in
polylogarithm order derivatives at explicit cyclotomic points. These
are convergent expansions, not finite period-basis assertions.
For example, \eqref{eq:Cm1} and the half-shift zeta formula give
\begin{align}
 C_{0,1}(\tfrac12,\tfrac34)
 ={}&\gamma_1-2\gamma\log2-\log^22\notag\\
 &+\sum_{k\ge1}\frac{(-1)^{k+1}(2^{k+1}-1)}{4^k k}\zeta(k+1).
 \label{eq:scalarquarter}
\end{align}
In the left side the cutoff subtraction is $\tfrac12\log^2N$.
The exact Gamma antiderivative \eqref{eq:C01primitive} then gives
\begin{align}
 \int_{1/2}^{3/4}C_{0,1}(\tfrac12,t)\dd t
  =\frac14C_{0,1}(\tfrac12,\tfrac34)
      -\log\Gamma(\tfrac34)+\tfrac12\log\pi
                           -\frac\gamma4-\frac12\log2.
 \label{eq:scalarquarterprimitive}
\end{align}
The constant itself has not been proved irreducible to a smaller
finite collection of special values.

\section{Source audit, reproducibility, and integration boundaries}
\label{sec:audit}
\subsection{What was actually inspected}
The targeted reading used the current manuscript README and editorial
ledger, together with the complete unequal-twist chapter and the
normalization portion of the preceding twisted chapter, pinned to
\pin. The incoming directory inventory was read, but the contents of
its binary ZIPs were not exhaustively retrieved or audited. A Library
search of related earlier reports corroborated the distinction between
previous shifted-argument correlation work and the present
unequal-frequency-shift problem. No inference of mathematical correctness
was made from an archive name or a successful earlier replay.

No new mathematical error was found in the two targeted canonical
chapters. In particular, the common-frequency-interval qualification of
the beta mixture, the unrestricted integer partial fractions, and the
separate status of the finite mixed-jet question are appropriate.
This report supplies an extension and a precise functional obstruction,
not a retraction of those results.

The ledger opening still describes a count of 439 textual source files,
whereas its later milestone and the README report 746. This is an
editorial synchronization issue; the full inventory was not independently
recounted here. The proposed correction is to make the opening use a
single authoritative inventory count or explicitly identify its historical
milestone, rather than silently treating both numbers as current.

\subsection{Mathematical safeguards for integration}
The following distinctions are essential. Neither the $k=0$ term of the central expansion nor the zero Fourier
coefficient may be deleted in the twisted calculus. The coefficient
$(0)_k$ must not be set to zero before its order derivatives are taken.
The integer partial-fraction theorem is not a holomorphic identity in
arbitrary orders and cannot simply be differentiated in those integer
labels. Finally, the contact term in \eqref{eq:Usymmetric}, the phase
second derivative in \eqref{eq:gammazero}, and the conductor-log term
in \eqref{eq:Stieltjesfilter} must be retained.

These are demonstrated hazards, not claims that the audited manuscript
made these errors. In the new multishift family, a purported joint
value at the zero-order intersection is specifically ruled out by
\eqref{eq:rayvalue}; a ray prescription must be stated.

\subsection{Executed checks and their evidentiary role}
The package includes \pathref{scripts/verify.py}, with separate
\texttt{exact}, \texttt{spectral}, \texttt{scalar}, and \texttt{gamma}
suites. A supplemental verifier, \pathref{scripts/verify_extended.py},
checks the general jet compiler and further analytic cases. The exact
checks contain 824 assertions in total, grouped as follows.
\begin{center}
\begin{tabular}{lr}
\toprule
Finite algebraic family & Assertions\\
\midrule
Pochhammer derivatives, including order zero &390\\
Multicentre coefficient recurrence &27\\
Symmetric mixed-log coefficients and odd cancellation &48\\
Gamma-convolution harmonic coefficients &25\\
Cutoff-jet Stirling coefficients &225\\
Mixed monodromy certificates &25\\
Single-centre mixed-monodromy cancellation &20\\
Integer partial fractions &16\\
Deliberately corrupted formulas rejected &3\\
General mixed-jet compiler against independent finite series &45\\
\midrule
Total &824\\
\bottomrule
\end{tabular}
\end{center}
The numerical suites make 38 spectral comparisons, 28 scalar
comparisons, four independent ordinary Gamma-integral comparisons, and
26 supplemental comparisons, for 96 numerical comparisons in total at
each of 65 and 85 working decimal digits.
The scalar comparison for $C_{m,n}$ uses the convergent difference sum
\eqref{eq:Cconvergent} with a separately computed Euler--Maclaurin
tail and a transformed bounded-interval integral; it does not reuse the
Hurwitz expansion of the mixed summand. The ordinary integral tests
use the finite log-Gamma representatives and compare them with
\eqref{eq:gammazero}--\eqref{eq:gammahalf}. Individual residuals and
the actual precision settings are recorded in the JSON reports.
The supplemental comparisons cover finite Fourier repairs across frequency
intervals, ordinary primitives of orders one through three, the pole-aware
Stieltjes root filter, the printed quarter-shift constant, and Cauchy
extraction of the first three spectral-ray derivatives.

A direct finite-difference call to the numerical polylogarithm routine
near order one was unstable at a root of unity. It was replaced by an
Euler-transformed defining series, with 150 extra working digits protecting
its finite differences. This computational issue is recorded separately
from the mathematical identities. The delivered root-filter diagnostics
use that independent convergent evaluator rather than the unstable call.

These checks test signs, finite coefficients, branches, zero modes,
normalizations, and concrete analytic instances. They do not replace
the proofs, establish a numerical independence theorem, or provide
interval-certified error bounds. The general distributional and
monodromy arguments are ordinary mathematical proofs, not
proof-assistant formalizations.

\subsection{Suggested placement}
The natural home is immediately after
\pathref{chapters/07-mixed-twist-identities.tex}, in the Hurwitz/Stieltjes
volume, with cross-references to the harmonic-number coefficient and
polylogarithm chapters. The package includes a short canonical addendum
with its own \texttt{utj:} labels, an integration map, a theorem-status
ledger, and replay/build instructions. The full report should remain as
a provenance source. None of these files has been pushed to the repository.

\section{Further research questions}
\label{sec:future}
\paragraph{1. Arithmetic specializations after functional nonclosure.}
For rational $a,b,c$ and rational $x$, identify combinations of
\eqref{eq:Esymmetric} or \eqref{eq:gammaconv} that collapse to a finite
polylogarithmic or Gamma basket. Theorem~\ref{thm:classification} does
not exclude these isolated identities. A first concrete target is a
finite evaluation, or a nontrivial finite relation, for the constants
in \eqref{eq:scalarquarter} and their higher logarithmic analogues.

\paragraph{2. Larger finite functional spaces.}
Determine the precise closure criterion when arbitrary complex orders,
polynomial coefficients in the argument, or finitely many spatial
translations are allowed. Each enlargement changes the frequency
space. The present theorem allows arbitrary finitely many centres,
but only integer order ladders and argument-independent coefficients.
An obstruction must be proved anew in the enlarged space.

\paragraph{3. Minimal mixed-centre logarithmic bases.}
Monodromy provides a support invariant recording which centres carry
logarithms. Construct a canonical finite reduction at fixed denominator
orders and logarithmic degrees in a basis that retains these supports.
The rational partial-fraction step should be separated from the
cross-centred logarithms rather than forcing them into a single-centre
basket which has just been proved insufficient.

\paragraph{4. Higher ordinary Gamma products.}
Apply the all-jet formula to three or more first primitives and seek
explicit symmetry reductions of their finite harmonic coefficients.
The multicentre theorem already supplies a convergent representation;
the remaining goal is compact identities, especially at rational
centres and selected circular arguments, not another existence proof.

\paragraph{5. Weighted spectra and higher Gamma functions.}
Replace the summand in $\mathcal Z$ by a polynomial multiplicity
$P(n)\prod_j(n+a_j)^{-\alpha_j}$. Compute the full ray-normalization
and first-derivative correction, comparing it with Barnes-type Gamma
functions. The cancellation leading to \eqref{eq:rayGamma} uses the
unweighted one-dimensional spectrum and should not be assumed to
survive unchanged.

\paragraph{6. Approach to an integral frequency shift.}
As a twist tends to an integer, isolate the exceptional Fourier mode
before taking order derivatives. Combine the finite-head formula
\eqref{eq:headrepair} with the canonical untwisted zero-mode subtraction
to obtain a uniform bridge between the unequal-twist and untwisted
Stieltjes calculi. The limits and normalization should be specified
jointly, rather than inferred from pointwise off-endpoint values.

\paragraph{7. Multivariable finite parts at intersecting divisors.}
The ray formula computes every transverse one-variable jet, but does
not select a canonical multivariable finite-part prescription. Determine
which coordinate changes preserve a chosen prescription for
\eqref{eq:multizeta}, and express the defects by its explicit residue
polynomials. This is distinct from assigning a joint value where the
ray limits demonstrably disagree.

\paragraph{8. Formal verification and independent analytic review.}
The finite coefficient recurrence, Stirling identities, and monodromy
finite differences are suitable first targets for proof-assistant
formalization. A separate analytic layer should formalize the
large-integer uniqueness lemma, Fourier pairing estimates, and
termwise differentiation. Interval-certified numerical variants may
then supplement, rather than substitute for, these proofs.

\section*{Conclusion}
\addcontentsline{toc}{section}{Conclusion}
Unequal twists exhibit an exact dichotomy. Integer rational multipliers
admit finite single-centre partial fractions, while genuine multishift
logarithmic decoration generally does not. A convergent central expansion
replaces the failed finite closure and gives explicit harmonic coefficients
at every order. Its primitives are ordinary Gamma convolution identities;
its unilateral counterpart gives polylogarithmic primitives, Hurwitz-jet
cutoff constants, and controlled spectral-ray finite parts. The distinction
between a finite functional identity, an infinite exact expansion, and an
isolated arithmetic value remains essential throughout.

\begin{thebibliography}{10}
\raggedright
\bibitem{repo-mixed}
ProveIt Contributors, \emph{Unequal twists through beta mixtures and finite
resolvent identities}, canonical manuscript source
\pathref{Analysis/Polylogarithms/docs/manuscript/chapters/07-mixed-twist-identities.tex},
revision \pin.
\href{https://github.com/VladimirReshetnikov/ProveIt/blob/b1df802799851c62a860570c3ac78ee31b9f5c04/Analysis/Polylogarithms/docs/manuscript/chapters/07-mixed-twist-identities.tex}{Pinned source on GitHub}.

\bibitem{repo-twisted}
ProveIt Contributors, \emph{Nonintegral twists: Lerch jets, contact terms,
and the half-twist}, canonical manuscript source
\pathref{Analysis/Polylogarithms/docs/manuscript/chapters/07-twisted-stieltjes.tex},
revision \pin. Targeted reading: family and normalization portions.
\href{https://github.com/VladimirReshetnikov/ProveIt/blob/b1df802799851c62a860570c3ac78ee31b9f5c04/Analysis/Polylogarithms/docs/manuscript/chapters/07-twisted-stieltjes.tex}{Pinned source on GitHub}.

\bibitem{repo-ledger}
ProveIt Contributors, \emph{Source reconciliation and mathematical
corrections}, \pathref{Analysis/Polylogarithms/docs/manuscript/EDITORIAL-LEDGER.md},
with the manuscript README; inspected 10 October 2026.
\href{https://github.com/VladimirReshetnikov/ProveIt/tree/main/Analysis/Polylogarithms/docs/manuscript}{Manuscript directory on GitHub}.

\bibitem{dlmf-zeta}
NIST Digital Library of Mathematical Functions, \S25.11,
\emph{Hurwitz Zeta Function}. In particular the defining series,
shift law, polygamma values, argument derivatives, and the log-Gamma
value of $\zeta'(0,a)$. \url{https://dlmf.nist.gov/25.11}.

\bibitem{dlmf-lerch}
NIST Digital Library of Mathematical Functions, \S25.14,
\emph{Lerch's Transcendent}. \url{https://dlmf.nist.gov/25.14}.

\bibitem{dlmf-gamma}
NIST Digital Library of Mathematical Functions, \S5.7,
\emph{Series Expansions of the Gamma and Psi Functions}.
\url{https://dlmf.nist.gov/5.7}.

\bibitem{dlmf-polylog}
NIST Digital Library of Mathematical Functions, \S25.12,
\emph{Polylogarithms}. \url{https://dlmf.nist.gov/25.12}.

\bibitem{dlmf-polygamma}
NIST Digital Library of Mathematical Functions, \S5.15,
\emph{Polygamma Functions}. \url{https://dlmf.nist.gov/5.15}.

\bibitem{lagarias-li}
J. C. Lagarias and W.-C. Winnie Li,
\emph{The Lerch Zeta Function III. Polylogarithms and Special Values},
Research in the Mathematical Sciences \textbf{3} (2016), article 2.
\url{https://doi.org/10.1186/s40687-015-0049-2}.
\end{thebibliography}
\end{document}
